\section{框架:线性解码结构、可模拟约化与单调性}\label{sec:framework}

\subsection{线性解码结构与噪声实例}

为了让同一套证明同时覆盖码容量噪声和电路级噪声(检测器误差模型),我们只用到\emph{线性性}。

\begin{definition}[线性解码结构]\label{def:structure}
一个线性解码结构是三元组 $\mathfrak D=(V,\sigma,L)$:$V$ 是有限维 $\F$-线性空间(\emph{故障空间}),
$\sigma:V\to\F^{r}$ 是\emph{检测器映射}(综合征),$L:V\to\F^{q}$ 是\emph{可观测量映射},均为线性映射。
\end{definition}

\begin{definition}[稳定子码的结构]\label{def:code-structure}
设 $S$ 为 $n$ 比特稳定子码的稳定子群($n-k$ 个独立生成元 $g_i$),泡利算符模相位记为 $V=\F^{2n}$,
辛内积 $\langle u,v\rangle=u_x\!\cdot v_z+u_z\!\cdot v_x$。取
\[
\sigma(E)=\big(\langle E,g_i\rangle\big)_{i=1}^{n-k},\qquad
L(E)=\big(\langle E,\bar X_j\rangle,\langle E,\bar Z_j\rangle\big)_{j=1}^{k},
\]
其中 $\bar X_j,\bar Z_j$ 是一组逻辑算符(辛基)。映射 $E\mapsto(\sigma(E),L(E))$ 满射,其核恰为 $S$。
电路级噪声:$V=\F^{\#\text{故障位置}}$,$\sigma$ 为检测器,$L$ 为可观测量翻转,线性性来自
Clifford 电路中泡利故障的线性传播。
\end{definition}

\begin{definition}[噪声实例与陪集区分问题]\label{def:instance}
噪声实例是随机对 $(E,F)$:$E\in V$ 是故障向量,$F\in\mathfrak F$($\mathfrak F$ 有限)是经典\emph{边信息}
(擦除标记、泄漏检测标记、事件标记等),联合分布 $P_{EF}$ 为解码器所知。解码器的观测为
$Y=(\sigma(E),F)$,任务是猜测 $L(E)$。解码器是(可随机化的)映射 $\hat g(Y)$。定义
\[
\err(P):=\min_{\hat g}\Prob[\hat g(Y)\neq L(E)]
=1-\sum_{y}\max_{\ell}\Prob[L=\ell,\,Y=y].
\]
最小值由最大似然(陪集)解码达到,$\err(P)=\pL$ 即最大似然下的逻辑失败概率。
\end{definition}

\begin{lemma}[解码即标签猜测]\label{lem:decoding}
对稳定子码(定义~\ref{def:code-structure}),设解码器看到 $y=(s,f)$ 后施加恢复算符 $R$,$\sigma(R)=s$。
则恢复成功(即 $E+R\in S$,逻辑恒等)当且仅当 $L(R)=L(E)$;且对每个 $(s,\ell)$ 存在 $R$ 使
$\sigma(R)=s,L(R)=\ell$。因此恢复失败概率的最小值就是 $\err(P)$。\hfill\st{PROVED}
\end{lemma}
\begin{proof}
$\sigma(E+R)=0$,故 $E+R$ 属于归一化子,且 $E+R\in S\iff L(E+R)=0\iff L(R)=L(E)$
(因 $(\sigma,L)$ 的联合核为 $S$)。$(\sigma,L)$ 满射给出第二句。
\end{proof}

\begin{remark}[标签选取无关]
把 $L$ 换成 $L'(E)=\pi_{\sigma(E)}(L(E))$,其中 $\pi_s$ 是依赖综合征的标签置换(例如换一个参考提升 $\rho(s)$),
$\err$ 不变:对固定 $y$ 取最大的 $\ell$ 的值不受置换影响。数值上我们直接使用原始标签
$(\langle E,\bar X\rangle,\langle E,\bar Z\rangle)$(见 \texttt{theory/checks})。
\end{remark}

对一族噪声模型 $\Nn=(P_d)_{d}$(码距 $d$ 为指标),记
$\alpha_-(\Nn)=\liminf_{d\to\infty}-\tfrac1d\ln\err(P_d)$,$\alpha_+(\Nn)=\limsup_{d\to\infty}-\tfrac1d\ln\err(P_d)$;
二者相等时记为压制指数 $\alpha$。

\subsection{可模拟约化}

\begin{definition}[可模拟约化]\label{def:reduction}
设 $(L,Y)$ 与 $(L',Y')$ 是取值于有限集的随机对。一个\emph{可模拟约化}
$(L,Y)\rightsquigarrow(L',Y')$ 是三元组 $(\kappa,g,h)$:
$\kappa(\cdot\mid y)$ 是从 $Y$ 到\emph{处理器随机性} $\xi$ 的马尔可夫核;$g(y,\xi)$ 与 $h(\ell,y,\xi)$ 为函数,满足

\smallskip
(i) 对每个 $(y,\xi)$,$\ell\mapsto h(\ell,y,\xi)$ 是标签集上的双射;

(ii) 令 $\xi\sim\kappa(\cdot\mid Y)$(在给定 $Y$ 时与 $L$ 条件独立),则
$\big(h(L,Y,\xi),\,g(Y,\xi)\big)$ 与 $(L',Y')$ 同分布。
\end{definition}

\begin{theorem}[单调性]\label{thm:monotone}
若 $(L,Y)\rightsquigarrow(L',Y')$,则 $\err(L\mid Y)\le\err(L'\mid Y')$。\hfill\st{PROVED}
\end{theorem}
\begin{proof}
设 $\hat g'$ 是 $(L',Y')$ 上的任一解码器。构造 $(L,Y)$ 上的随机化解码器 $\hat g$:
观察 $y$,抽 $\xi\sim\kappa(\cdot\mid y)$,输出 $h^{-1}(\hat g'(g(y,\xi));y,\xi)$,其中 $h^{-1}$ 对标签变量取逆(由 (i) 存在)。
则
\[
\{\hat g(Y)=L\}=\{\hat g'(g(Y,\xi))=h(L,Y,\xi)\}=\{\hat g'(Y'_{\rm sim})=L'_{\rm sim}\},
\]
其中 $(L'_{\rm sim},Y'_{\rm sim})=(h(L,Y,\xi),g(Y,\xi))$ 由 (ii) 与 $(L',Y')$ 同分布。故
$\Prob[\hat g(Y)\neq L]=\Prob[\hat g'(Y')\neq L']$。随机化解码器的错误率是确定性解码器错误率的平均,
所以存在确定性解码器的错误率不超过它。对 $\hat g'$ 取最优即得结论。
\end{proof}

\begin{remark}[数据处理不等式的形式]
$1-\err(L\mid Y)=2^{-H_{\min}(L\mid Y)}$($H_{\min}$ 为条件最小熵)。当 $h=\mathrm{id}$
(只处理观测)时,定理~\ref{thm:monotone} 就是 $H_{\min}(L\mid Y)\le H_{\min}(L\mid g(Y,\xi))$,
即条件最小熵(Sibson--R\'enyi $\infty$ 阶互信息)的数据处理不等式。其余 R\'enyi/Hellinger 型泛函同样单调,
见命题~\ref{prop:hellinger}。
\end{remark}

\subsection{三类自由操作}

下面三类操作都是定义~\ref{def:reduction} 的特例,从而由定理~\ref{thm:monotone} 得到单调性。
设 $(E,F)$ 是线性解码结构 $\mathfrak D$ 上的噪声实例。

\begin{description}[leftmargin=2.2em,style=nextline]
\item[(F1) 综合征粗粒化]
$Y'=g(Y,\xi)$,$h=\mathrm{id}$,$g$ 为任意(可随机)函数,例如合并检测器、取子集。\\
\emph{特例}:丢弃标记(\textbf{(F1$'$)}):$F'=\varphi(F,\xi)$,仍是噪声实例 $(E,F')$。
\item[(F2) 叠加可经典采样的独立噪声]
设 $(E',F'')$ 是与 $(E,F)$ 独立、分布已知的随机对。新实例为
$(E+E',\,(F,F''))$(或经 (F1$'$) 进一步粗粒化的版本)。
取 $\xi=(E',F'')$,$g(y,\xi)=(\sigma(E')+s,\,f,\,F'')$,$h(\ell,y,\xi)=\ell+L(E')$。
\item[(F2$'$) 依赖已观测标记的独立噪声]
$\xi\sim\kappa(\cdot\mid y)$ 可依赖观测 $y=(s,f)$;例如只对未被标记的比特叠加去极化噪声。
$h(\ell,y,\xi)=\ell+L(N)$,$N$ 为处理器自己抽出的叠加故障。
\end{description}

\begin{proposition}[(F1)(F2)(F2$'$) 为可模拟约化]\label{prop:free-ops}
设 $(E,F)$ 为噪声实例,按 (F2) 或 (F2$'$) 得到 $(E_{\rm new},F_{\rm new})$,其中叠加的故障 $N$ 满足:
给定 $F$ 时 $N$ 与 $E$ 条件独立(特别地,(F2) 中 $N=E'$ 与 $(E,F)$ 独立)。则
\[
\err(E,F)\;\le\;\err(E+N,\,F_{\rm new}),
\]
其中 $F_{\rm new}$ 是 $(F,\text{附加标记})$ 的任一(随机)函数。特别地,叠加噪声永不降低 $\err$。
\hfill\st{PROVED}
\end{proposition}
\begin{proof}
新问题的标签为 $L(E+N)=L(E)+L(N)$(线性性),观测为 $(\sigma(E)+\sigma(N),F_{\rm new})$。
处理器由 $(\sigma(E),F)$ 与自己抽取的 $(N,\text{附加标记})$ 算出新观测,并用已知的 $L(N)$ 做标签平移。
这是定义~\ref{def:reduction} 的约化,其中 (i) 因平移是双射,(ii) 因 $N$ 条件独立于 $E$ 且分布已知
给出正确联合分布。再应用定理~\ref{thm:monotone}。
\end{proof}

\begin{corollary}[指数的单调性]\label{cor:alpha}
若对每个 $d$,$P'_d$ 由 $P_d$ 经(同类型的)自由操作得到,则
$\alpha_\pm(\Nn')\le\alpha_\pm(\Nn)$。\hfill\st{PROVED}
\end{corollary}

\begin{remark}[独立性不可去掉]\label{rem:independence}
取 $N=E$(叠加的噪声是已有故障的拷贝):$E+N=0$,$\err=0<\err(E,F)$。所以
处理器必须能在\emph{不知道隐变量 $E$} 的情况下抽出 $N$——这正是“叠加\emph{独立}噪声”的含义。
第~\ref{sec:repr} 节据此划定模型类的边界。
\end{remark}

\begin{remark}[只对最优解码器成立]\label{rem:ml-only}
定理~\ref{thm:monotone} 刻画的是 Bayes 最优(最大似然)解码。对次优解码器(如固定权重的 MWPM),
叠加噪声并不保证提高其失败率的单调性,也没有现成的数据处理不等式;因此单调性的\emph{数值}检验必须使用精确 ML
(见 \texttt{theory/checks/exact\_small\_codes.py}),不能用 MWPM 结果代替。
反过来,由最优性始终有 $\err\le p_{\rm L}^{\rm MWPM}$,即 $\alpha_{\rm ML}\ge\alpha_{\rm MWPM}$。
\end{remark}
